\documentclass[addpoints,fleqn,11pt]{exam}

% -----------------------------------------------------
% EDITAR - Información del Curso, Unidad Académica y Documento
% -----------------------------------------------------
\newcommand{\unidadAcademica}{Facultad de Matemáticas} % EDITAR - Unidad Académica/Organización (ej. Facultad de Matemáticas)
\newcommand{\logoUnidadAcademica}{delegados2020.png}   % EDITAR - Logo Unidad Académica/Organización (ej. delegados2020.png)
\newcommand{\codigoDelCurso}{MAT1610}                  % EDITAR - Código del Curso (ej. MAT1203)
\newcommand{\elCurso}{Cálculo I}                       % EDITAR - Título del Curso (ej. Cálculo I)
\newcommand{\fecha}{Mayo 2020}                         % EDITAR - Fecha del Documento (ej. Mayo 2020)
\newcommand{\nombreDocumento}{I1}                      % EDITAR -  nombre del documento. (Libre, ej. Interrogación 1 o Guía de Estudio)
\newcommand{\tipoDocumento}{interrogación}             % EDITAR - tipo del documento. (Puede ser interrogación, control, estudiatón o guía)

\newcommand{\descripcionPortada}{%                     % EDITAR - Descripción que aparece como primer párrafo de la portada.
        Esta guía fue hecha con {\emoji \symbol{"2764"}} por novatos para novatos, para la estudiatón de cálculo el 13 de mayo de 2020, específicamente en preparación a la I2. Incluso si no pudiste estar en la estudiatón, esperamos que esta guía te sea útil para practicar y prepararte.\\[\baselineskip]
    }
\newcommand{\instruccionesPortada}                     % EDITAR - Instrucciones que aparecen al final de la portada.
{%
    \textbf{Si tienes alguna duda, primero intenta llevarlo al grupo en el que estás.} Si nadie puede resolver su duda, puedes solicitar ayuda desde Zoom para que vaya un tutor a ayudar. Si todos en tu grupo terminan la guía, o están muy pegados con una pregunta de desarrollo, puedes pedir la pauta de esta guía.
}

\newcommand{\hojaDeFormulas}{Si}                       % EDITAR - ¿Incluir hoja de fórmulas? (Si o no)
\newcommand{\tieneSecciones}{No}                       % EDITAR - ¿El documento tiene secciones? (Si o no)
                                                       %        Si: Usa \iniciarSeccion{} para empezar secciones.
                                                       %        No: Comenta todos los \iniciarSeccion{}
\newcommand{\conRespuestas}{No}                        % EDITAR - Incluir soluciones (Si o no)
\newcommand{\soportaEmojis}{Si}                        % EDITAR - ¿El documento soporta emojis? (Si o no, requiere lualatex-dev)
\newcommand{\soportaCodigo}{Si}                        % EDITAR - ¿El documento soporta bloques de código? (Si o no, requiere minted, pygments de Python)
\newcommand{\usaOverleaf}{No}                          % EDITAR - Si se está usando Overleaf como editor, es necesario configurar esto como "Si", corrigiendo un error de compilación.

% Para evaluaciones (solo aplican si es que \tipoDocumento es control o interrogación)
\newcommand{\mostrarPuntos}{Si}                        % EDITAR - Mostrar puntos de cada pregunta. (Si o no)
\newcommand{\libroAC}{Cerrado}                         % EDITAR - Tipo de Examen (Abierto o Cerrado)
\newcommand{\preguntasPorResponder}{Todas}             % EDITAR - Número de preguntas a responder del total, (ej. Todas o 3)
\newcommand{\conCalculadoras}{No}                      % EDITAR - ¿Se permiten calculadoras? (Si o no)
\newcommand{\respuestaPorSeparado}{Si}                 % EDITAR - Se responde las preguntas en un documento separado? (Si o no)

% -----------------------------------------------------
% INICIO DEL DOCUMENTO
% -----------------------------------------------------
% !TeX root = ./documento.tex
% -------------------------------------------------------------
% Paquetes
% -------------------------------------------------------------
\usepackage{microtype}
\usepackage{commath}                                  % Configura el espacio tipográfico por razones estéticas
\usepackage{cmbright}                                 % Tipografía (Computer Modern Bright)
\usepackage[colorlinks = true,
    allcolors = black]{hyperref}                      % Vínculos
\usepackage{subfigure}                                % Entorno de multiples figuras
\usepackage{amsmath}                                  % Matemáticas estándar
\usepackage{amssymb}                                  % Símbolos matemáticos estándar
\usepackage{graphicx}                                 % Para incluir imágenes
\usepackage{booktabs}                                 % Tablas profesionales
\usepackage[spanish, es-noshorthands]{babel}          % Definición del lenguaje

% Deshabilitado: Tamaño A4
% \usepackage[a4paper,                                % Tamaño del papel
%     top = 25mm,                                     % Margen superior
%     bottom = 25mm,                                  % Margen inferior
%     left = 25mm,                                    % Margen izquierdo
%     right = 25mm]{geometry}                         % Margen derecho

\usepackage[dvipsnames]{xcolor}                       % Colores avanzados
\usepackage{soul}                                     % Resaltar texto y otras cosas
\usepackage{caption}                                  % Descripciones en objetos flotantes
\usepackage{ifthen}                                   % Condicionales if/then
\usepackage{fmtcount}                                 % Convertir enteros a palabras
\usepackage[per-mode = symbol]{siunitx}	              % Unidades del S.I.U
\usepackage[version=4]{mhchem}                        % Símbolos químicos
\usepackage{totcount}                                 % Mostrar contador de puntaje al princpio
\usepackage{xstring}                                  % Para manipular strings
\usepackage{multicol}                                 % Entorno de multiples columnas (FS)
\usepackage{mfirstuc}                                 % Agrega comandos para mayúsculas
\usepackage[no-math]{fontspec}                        % Manejo tipográfico manual
\usepackage{wasysym}                                  % Símbolos para preguntas de alternativas
\usepackage{etoolbox}                                 % if/else modernos
\usepackage[outputdir=out]{minted}                    % Soporte para código

% -----------------------------------------------------
% Definiciones del Documento
% -----------------------------------------------------

\renewcommand{\partlabel}{\bfseries{(\thepartno)}}     % Números de parte negrita
\renewcommand{\subpartlabel}{\bfseries{\thesubpart{}}} % Números de parte negrita
\pagestyle{headandfoot}                                % Header y Footer de la solución
\header{}{}{}                                          % Header
\footer{}{}{\codigoDelCurso{}                          % Footer
    \elCurso{} --- \fecha{}
    (\nombreDocumento)}
\graphicspath{{imagenes/}}                             % Directorio con las imágenes
\definecolor{myRed}{RGB}{184,0,1}                      % Nuevo gris
\definecolor{myGrey}{RGB}{211,211,211}                 % Nuevo rojo
\checkedchar{\large \(\CIRCLE\)}
\setlength{\fboxrule}{0pt}

% -----------------------------------------------------
% Versiones
% -----------------------------------------------------
\newcommand{\theVersion}{1.3}
\newcommand{\theVersionDate}{24-May-2020}

% -----------------------------------------------------
% Nuevos Comandos
% -----------------------------------------------------

\newtotcounter{contadorDePuntos}                      % Contador del número de puntos por parte
\newtotcounter{numeroDeSecciones}                     % Contador de número de secciones
\newtotcounter{numeroDePreguntas}                     % Contador de número de preguntas
\DeclareMathOperator{\arcsec}{arcsec}                 % Funciones Trigonométricas Inversas
\DeclareMathOperator{\arccot}{arccot}
\DeclareMathOperator{\arccsc}{arccsc}

% -----------------------------------------------------
% Lógica Interna
% -----------------------------------------------------

% Thanks Phelype Oleinik for his awesome answer on StackExchange suggesting this. https://tex.stackexchange.com/users/134574/phelype-oleinik
\newrobustcmd\sanitizadorSiNo[1]{%
    \MakeLowercase{\gdef\noexpand\internalCondition{#1}}%
    \ifboolexpr{%
        test {\expandafter\ifstrequal\expandafter{\internalCondition}{yes}}
        or
        test {\expandafter\ifstrequal\expandafter{\internalCondition}{yas}}
        or
        test {\expandafter\ifstrequal\expandafter{\internalCondition}{si}}
        or
        test {\expandafter\ifstrequal\expandafter{\internalCondition}{sí}}
        or
        test {\expandafter\ifstrequal\expandafter{\internalCondition}{s\IeC{\'\i}}}
        or
        test {\expandafter\ifstrequal\expandafter{\internalCondition}{s\IeC{\'i}}} % For SÍ
    }}

\newrobustcmd\sanitizadorPorResponder[1]{%
    \MakeLowercase{\gdef\noexpand\internalCondition{#1}}%
    \ifboolexpr{%
        test {\expandafter\ifstrequal\expandafter{\internalCondition}{todas}}
        or
        test {\expandafter\ifstrequal\expandafter{\internalCondition}{todo}}
        or
        test {\expandafter\ifstrequal\expandafter{\internalCondition}{toda}}
        or
        test {\expandafter\ifstrequal\expandafter{\internalCondition}{all}}
    }}
\newrobustcmd\sanitizadorLibroAC[1]{%
    \MakeLowercase{\gdef\noexpand\internalCondition{#1}}%
    \ifboolexpr{%
        test {\expandafter\ifstrequal\expandafter{\internalCondition}{abierto}}
        or
        test {\expandafter\ifstrequal\expandafter{\internalCondition}{apuntes}}
    }}
\newrobustcmd\sanitizadorEsEvaluacion[1]{%
    \MakeLowercase{\gdef\noexpand\internalCondition{#1}}%
    \ifboolexpr{%
        test {\expandafter\ifstrequal\expandafter{\internalCondition}{interrogación}}
        or
        test {\expandafter\ifstrequal\expandafter{\internalCondition}{interrogacion}}
        or
        test {\expandafter\ifstrequal\expandafter{\internalCondition}{control}}
        or
        test {\expandafter\ifstrequal\expandafter{\internalCondition}{evaluación}}
        or
        test {\expandafter\ifstrequal\expandafter{\internalCondition}{evaluacion}}
    }}

\newrobustcmd\siHojaDeFormulas{%
    \expandafter\sanitizadorSiNo
    \expandafter{\hojaDeFormulas}}%
\newrobustcmd\siTieneSecciones{%
    \expandafter\sanitizadorSiNo
    \expandafter{\tieneSecciones}}%
\newrobustcmd\siConRespuestas{%
    \expandafter\sanitizadorSiNo
    \expandafter{\conRespuestas}}%
\newrobustcmd\siSoportaEmojis{%
    \expandafter\sanitizadorSiNo
    \expandafter{\soportaEmojis}}%
\newrobustcmd\siSoportaCodigo{%
    \expandafter\sanitizadorSiNo
    \expandafter{\soportaCodigo}}%
\newrobustcmd\siMostrarPuntos{%
    \expandafter\sanitizadorSiNo
    \expandafter{\mostrarPuntos}}%
\newrobustcmd\siConCalculadoras{%
    \expandafter\sanitizadorSiNo
    \expandafter{\conCalculadoras}}%
\newrobustcmd\siRespuestaPorSeparado{%
    \expandafter\sanitizadorSiNo
    \expandafter{\respuestaPorSeparado}}%
\newrobustcmd\siResponderTodo{%
    \expandafter\sanitizadorPorResponder
    \expandafter{\preguntasPorResponder}}%
\newrobustcmd\siEsEvaluacion{%
    \expandafter\sanitizadorEsEvaluacion
    \expandafter{\tipoDocumento}}%
\newrobustcmd\siEsLibroAbierto{%
    \expandafter\sanitizadorLibroAC
    \expandafter{\libroAC}}%
\ExplSyntaxOff

% Define comandos y entornos en español
\newcommand{\pregunta}{\question}
\newcommand{\mostrarpuntaje}{\droppoints}
\newenvironment{solucion}{\begin{solution}}{\end{solution}}
\newenvironment{alternativas}{\begin{checkboxes}}{\end{checkboxes}}
\newcommand{\alternativa}{\choice}
\newcommand{\alternativaCorrecta}{\CorrectChoice}



% Maneja familia del documento
\siEsEvaluacion{\newcommand{\familiaDocumento}{evaluación}}{\newcommand{\familiaDocumento}{guía}}

% Maneja el soporte de emojis
\siSoportaEmojis{
    \newfontfamily{\emoji}{emoji.ttf}[Renderer=Harfbuzz]}  %  if "Enunciados con emojis" - ΠCAi
{}

% Maneja el soporte de bloques de código con minted
\siSoportaCodigo{
    \usepackage{minted}}
{}

% -----------------------------------------------------
% Estructura del Documento - Preguntas y secciones
% -----------------------------------------------------

\siTieneSecciones
{\newcommand{\sectionStart}{                              % Iniciar sección
        \setcounter{question}{0}                          % Reiniciar contador preguntas
        \setcounter{figure}{0}                            % Iniciar contador de figuras
        \setcounter{table}{0}                             % Iniciar contador de tablas
        \qformat{\ifthenelse{\equal{\thequestion}{1}}
            {\fbox{\parbox{16cm}{\textbf{\large SECCIÓN \Alph{numeroDeSecciones} \hfill\vrule depth 0.5em width 0pt \\
                            Question \Alph{numeroDeSecciones}\thequestion}} \hfill\vrule depth 1.5em width 0pt}}
            {\fbox{\parbox{16cm}{\large \textbf{Pregunta \Alph{numeroDeSecciones}\thequestion}} \hfill\vrule depth 0.8em width 0pt}}
        }
        \stepcounter{numeroDeSecciones}               % Incrementar contador
    }}

% Si no hay secciones
{   \setcounter{question}{0}                          % Reinicializar contador de preguntas
    \setcounter{figure}{0}                            % Iniciar contador de figuras
    \setcounter{table}{0}                             % Iniciar contador de tablas
    \qformat{
        {\fbox{\parbox{16cm}{\large \textbf{Pregunta \thequestion}} \hfill\vrule depth 0.8em width 0pt}}
    }
}

% Invocador misterioso de preguntas
\newcommand{\traerPregunta}[2]{                       % Nueva pregunta
    \cleardoublepage \input{#1}                       % Iniciar pregunta en nueva página
    \addtocounter{contadorDePuntos}{#2}               % Incrementar puntos
    \addtocounter{numeroDePreguntas}{1}               % Incrementar preguntas
    \stepcounter{figure}                              % Incrementar figuras
    \stepcounter{table}                               % Incrementar tablas
}

% Define la hoja de fórmulas
\newcommand{\traerHojaDeFormulas}[1]{
    \ifthenelse{
        \equal{\hojaDeFormulas}{Yes} \OR{} \equal{\hojaDeFormulas}{Y} \OR{} \equal{\hojaDeFormulas}{yes} \OR{} \equal{\hojaDeFormulas}{y} \OR{} \equal{\hojaDeFormulas}{YES} \OR{} \equal{\hojaDeFormulas}{Si} \OR{} \equal{\hojaDeFormulas}{Sí} \OR{} \equal{\hojaDeFormulas}{Verdadero} \OR{} \equal{\hojaDeFormulas}{si} \OR{} \equal{\hojaDeFormulas}{sí} \OR{} \equal{\hojaDeFormulas}{verdadero} \OR{} \equal{\hojaDeFormulas}{SÍ} \OR{} \equal{\hojaDeFormulas}{SI} \OR{} \equal{\hojaDeFormulas}{VERDADERO}
    }
    {\cleardoublepage \input{#1}}{}
}

% Cabecera de la hoja de fórmulas
\newcommand{\headingFormulaSheet}{
    \pagestyle{empty}
    \setlength{\textwidth}{180mm}
    \setlength{\textheight}{265mm}
    \setlength{\topmargin}{-20mm}
    \setlength\oddsidemargin{-12mm}
    \setlength\evensidemargin{-12mm}
    \setlength{\columnsep}{10mm}
    \setlength{\columnseprule}{0.4pt}
    \setlength{\jot}{5pt}
    \setlength\parindent{0pt}
    %
    \twocolumn
    \fcolorbox{white}{myGrey}{\begin{minipage}{20em}
            \centering \vspace{4mm} {\Large \textbf{Hoja de Fórmulas}} \\
            \elCurso\ (\codigoDelCurso) \\
            {\small \fecha{} (\tipoDocumento)} \vspace{3mm}
        \end{minipage}}
}

% -----------------------------------------------------
% Títulos y formatos
% -----------------------------------------------------

%\qformat{\textbf{\large Question \thequestion} \hfill\vrule depth 1.5em width 0pt}

\pointsinrightmargin{}                                % Puntos en el margen derecho
\siMostrarPuntos{\pointformat{\bfseries(\themarginpoints)}}{\pointformat{}}

% Formato del título de la pregunta - solo las numeradas y no en partes
\pointname{}                                          % No title for marks
\pointsdroppedatright{}                               % Puntaje por pregunta
\framedsolutions{}                                    % Marco para soluciones
\addto\captionsenglish{                               % Formato de etiquetas de figura
    %    \renewcommand\figurename{Figure \thequestion .}}
    \renewcommand\figurename{Figura}}
\addto\captionsenglish{                               % Etiquetas de tabla
    %    \renewcommand\tablename{Table \thequestion .}}
    \renewcommand\tablename{Tabla}}
\makeatletter
\def\fnum@figure{\figurename~\Alph{numeroDeSecciones}\thefigure}
\def\fnum@table{\tablename~\Alph{numeroDeSecciones}\thetable}
\makeatother
\renewcommand{\solutiontitle}{
    \noindent\textbf{Solución:}\par\noindent}  % Renombrar soluciones

% -----------------------------------------------------
% ¿Silenciar advertencias? Obvio!
% -----------------------------------------------------

% Desactiva errores de overfull y underfull.
\hfuzz=\maxdimen
\hbadness=10000                                           % Definiciones de estilo y estructura
\begin{document}                                       % Inicio del contenido

% -----------------------------------------------------
% Portada del Documento (automatizada, revisar documento.tex)
% -----------------------------------------------------
\begin{coverpages}
    \begin{flushleft}

        \includegraphics[height=45mm]{puc_comunitario.png}
        \includegraphics[height=45mm]{delegados2020.png} \\ \vspace{10mm}

        \textcolor{MidnightBlue}{\textbf{\Large \codigoDelCurso{}: \elCurso{} --- \fecha{} (\nombreDocumento)}} \\[\baselineskip]

        \descripcionPortada{}

        \ifthenelse{
            \equal{\hojaDeFormulas}{Yes} \OR{} \equal{\hojaDeFormulas}{Y} \OR{} \equal{\hojaDeFormulas}{yes} \OR{} \equal{\hojaDeFormulas}{y} \OR{} \equal{\hojaDeFormulas}{YES} \OR{} \equal{\hojaDeFormulas}{Si} \OR{} \equal{\hojaDeFormulas}{Sí} \OR{} \equal{\hojaDeFormulas}{Verdadero} \OR{} \equal{\hojaDeFormulas}{si} \OR{} \equal{\hojaDeFormulas}{sí} \OR{} \equal{\hojaDeFormulas}{verdadero} \OR{} \equal{\hojaDeFormulas}{SÍ} \OR{} \equal{\hojaDeFormulas}{SI} \OR{} \equal{\hojaDeFormulas}{VERDADERO}
        }
        {Este documento \textbf{incluye una hoja de fórmulas} hacia el final. Puedes usarlo como referencia.\\[\baselineskip]} % if

        \ifthenelse{\equal{\familiaDocumento}{evaluación}}{
            \ifthenelse{
                \equal{\preguntasPorResponder}{Todas} \OR
                \equal{\preguntasPorResponder}{todas}
            }
            {Para esta evaluación \textbf{debes intentar responder todas las preguntas.} \\[\baselineskip]} % if
            {Para esta evaluación, debes responder \textbf{solo} \NUMBERstringnum{\total{numeroDePreguntas}} preguntas. \\[\baselineskip]} % else
            \ifthenelse{
                \equal{\conCalculadoras}{Si} \OR \equal{\conCalculadoras}{Sí} \OR{} \equal{\conCalculadoras}{Verdadero} \OR{} \equal{\conCalculadoras}{si} \OR{} \equal{\conCalculadoras}{sí} \OR{} \equal{\conCalculadoras}{verdadero} \OR{} \equal{\conCalculadoras}{SÍ} \OR{} \equal{\conCalculadoras}{SI} \OR{} \equal{\conCalculadoras}{VERDADERO}
            }{
                El uso de calculadoras científicas autorizadas está \textbf{permitido} en esta evaluación.\\[\baselineskip]
            } % if
            {
                El uso de calculadoras de cualquier tipo está \textbf{prohibido} durante esta evaluación.\\[\baselineskip]
            } % else
            \ifthenelse{
                \equal{\libroAC}{Si} \OR \equal{\libroAC}{Sí} \OR{} \equal{\libroAC}{Verdadero} \OR{} \equal{\libroAC}{si} \OR{} \equal{\libroAC}{sí} \OR{} \equal{\libroAC}{verdadero} \OR{} \equal{\libroAC}{SÍ} \OR{} \equal{\libroAC}{SI} \OR{} \equal{\libroAC}{VERDADERO}
            }{
                Esta es una evaluación a \textbf{libro abierto}, por lo que puedes utilizar apuntes o libros acorde con el código de honor. El plagio se encuentra prohibido.\\[\baselineskip]
            }
            {
                Esta es una evaluación a \textbf{libro cerrado}, por lo que tienes \textbf{prohibido} el uso de cualquier clase de apuntes, libros o hojas de referencia, al menos que una sea incluida con el documento.\\[\baselineskip]
            }
            \ifthenelse{
                \equal{\respuestaPorSeparado}{Si} \OR \equal{\respuestaPorSeparado}{Sí} \OR{} \equal{\respuestaPorSeparado}{Verdadero} \OR{} \equal{\respuestaPorSeparado}{si} \OR{} \equal{\respuestaPorSeparado}{sí} \OR{} \equal{\respuestaPorSeparado}{verdadero} \OR{} \equal{\respuestaPorSeparado}{SÍ} \OR{} \equal{\respuestaPorSeparado}{SI} \OR{} \equal{\respuestaPorSeparado}{VERDADERO}
            }
            {
                Para responder esta evaluación, \textbf{debes escribir tus respuestas por separado}, en una hoja o documento que se te debería entregar. No es aceptable entregar las respuestas en esta hoja.\\[\baselineskip]
            }
            {
                Para responder esta evaluación, \textbf{debes escribir tus respuestas en este documento}. No es aceptable entregar las respuestas por separado, salvo que sea explícitamente solicitado.\\[\baselineskip]
            }
        }{}

        \vfill
        % Advertencia de respuestas
        \ifthenelse{\equal{\conRespuestas}{Si} \OR{} \equal{\conRespuestas}{Sí} \OR{} \equal{\conRespuestas}{si} \OR{} \equal{\conRespuestas}{sí} \OR{} \equal{\conRespuestas}{SÍ} \OR{} \equal{\conRespuestas}{SI}}{\textcolor{myRed}{\LARGE \textbf{Advertencia:} Esta copia incluye respuestas.} \\}{}
        \vfill

        \textcolor{MidnightBlue}{\textbf{\large Instrucciones Adicionales}}\\

        Todos en el Zoom van a ser divididos en grupos de \textbf{5 a 8 personas} al iniciar la estudiatón.\\ Una vez que seas asignado a un grupo, \textbf{tendrás que presentarte antes de empezar} y después podrán comenzar a hacer en orden cada pregunta.\\[\baselineskip]

        \textbf{Si tienen alguna duda, primero intenten llevarlo al grupo en el que están.} Si nadie puede resolver su duda, pueden solicitar ayuda desde Zoom para que vaya un tutor a ayudar. Si todos en tu grupo terminan la guía, o están muy pegados con una pregunta de desarrollo, pueden pedir la pauta de esta guía.

        \vspace{3mm} \textcolor{MidnightBlue}{\small Delegados Generacionales 2020 Ing. UC}
        %
    \end{flushleft}
\end{coverpages}
\newpage

% Impresión automática de respuestas, NO REMOVER %
\ifthenelse{
    \equal{\conRespuestas}{Yes} \OR{} \equal{\conRespuestas}{Y} \OR{} \equal{\conRespuestas}{yes} \OR{} \equal{\conRespuestas}{y} \OR{} \equal{\conRespuestas}{YES} \OR{} \equal{\conRespuestas}{Si} \OR{} \equal{\conRespuestas}{Sí} \OR{} \equal{\conRespuestas}{Verdadero} \OR{} \equal{\conRespuestas}{si} \OR{} \equal{\conRespuestas}{sí} \OR{} \equal{\conRespuestas}{verdadero} \OR{} \equal{\conRespuestas}{SÍ} \OR{} \equal{\conRespuestas}{SI} \OR{} \equal{\conRespuestas}{VERDADERO}
}
{\printanswers}{}                                    % Portada del Documento
\begin{questions}                                      % Inicio de las preguntas
    % -----------------------------------------------------
    % Iniciar preguntas
    % EDITAR - Copiar para añadir secciones o preguntas, una pregunta por archivo.
    % -----------------------------------------------------
    \traerPregunta{preguntas/alternativas.tex}{35}     % EDITAR - Archivo de la pregunta y su puntaje
    \traerPregunta{preguntas/doce.tex}{20}
    \traerPregunta{preguntas/trece.tex}{20}
    \traerPregunta{preguntas/catorce.tex}{20}
    % -----------------------------------------------------
    % Terminar preguntas
    % -----------------------------------------------------

\end{questions}                                        % Finalizar secciones y preguntas
\vfill \medskip\hfill\textbf{FIN DE LA GUÍA}           % Finalizar documento principal

\traerHojaDeFormulas{formulas.tex}                     % NO EDITAR. Se deshabilita automáticamente.

% -----------------------------------------------------
% Finalizar documento
% -----------------------------------------------------
\end{document}
